% \documentclass{article}
\documentclass[a4paper,11pt]{l3doc}
% \usepackage[margin=1cm]{geometry}
\usepackage[fontset=ubuntu]{ctex}
\usepackage{cube}
\usepackage[skins,documentation]{tcolorbox}
\usepackage{enumitem}
\usepackage{array,tabularx,colortbl}

% region 样式设置
\let\oldmarg=\marg
\renewcommand{\marg}[1]{\oldmarg{\kaishu#1}}
\let\oldoarg=\oarg
\renewcommand{\oarg}[1]{\oldoarg{\kaishu#1}}
\let\oldmeta=\meta
\renewcommand{\meta}[1]{\oldmeta{\kaishu#1}}
\renewcommand{\thesubsection}{\bullet}
\let\oldsubsection=\subsection
\renewcommand{\subsection}[1]{\oldsubsection{\texttt{\zihao{5}#1}}}
\DeclareDocumentCommand{\textsccolor}{ m }{\textcolor{Definition}{\textsc{\kaishu#1}}}
\DeclareDocumentCommand{\textsfcolor}{ m }{\textcolor{niceblue}{\textsf{\kaishu#1}}}
\DeclareDocumentCommand{\textttcolor}{ m }{\textcolor{Definition}{\texttt{\kaishu#1}}}
\DeclareRobustCommand \rcs {\textttcolor}
\DeclareRobustCommand \pkg {\textsfcolor}
\DeclareRobustCommand \env {\textttcolor}

\newtcolorbox{cubeExample}{%
  enhanced,
  colback=ExampleBack,
  colframe=ExampleFrame,
  before skip=10pt,
  after skip=10pt,
  boxsep=5pt,
}
% endregion

% region 自定义命令
\ProvideDocumentCommand{\red}{m}{\textcolor{red}{#1}}
\ProvideDocumentCommand{\blue}{m}{\textcolor{blue}{#1}}
\ProvideDocumentCommand{\orange}{m}{\textcolor{orange}{#1}}
\ProvideDocumentCommand{\green}{m}{\textcolor{nicegreen}{#1}}
% endregion

% region 颜色设置
\definecolor{spot}{HTML}{003399}
\definecolor{code}{HTML}{a25e26}
\definecolor{verb}{HTML}{007f00}
\let\OriginalMF\MacroFont
\def\MacroFont{\color{spot}\OriginalMF}
\makeatletter
\pretocmd{\macrocode}{\def\macro@font{\color{code}\OriginalMF}\def\MacroFont{\color{code}\OriginalMF}}{}{}
\makeatother
\AtBeginEnvironment{quote}{\color{verb}}
\AtBeginEnvironment{flushleft}{\tt\color{verb}}
\AtBeginEnvironment{verbatim}{\color{verb}\def\MacroFont{\color{verb}\OriginalMF}}
\ExplSyntaxOn
\pretocmd{\__codedoc_typeset_functions:}{\color{spot}\arrayrulecolor{spot}}{}{}
\patchcmd{\__codedoc_typeset_aux:n}{\color[gray]{0.5}}{\color{spot}}{}{}
\ExplSyntaxOff
\hypersetup{%
  colorlinks=true,
  linkcolor=spot,
  urlcolor=spot,
  citecolor=spot,
  bookmarksopen=false,
  bookmarksnumbered=false,
  plainpages=false}
% endregion
\begin{document}

\ttfamily % 设置全文英文字体为等宽字体

% region 标题
\title{{\raisebox{-1ex}{\CubeLogo[0.5]}}\pkg{cube}：魔方绘制宏包}
\author{张泓知}
\date{\today}
\maketitle
% endregion

% region 简介
\section{简介}
\pkg{cube} 宏包用于绘制魔方相关图形。它提供了简单的命令来表示三阶魔方的各个面、块和颜色，使得在 \LaTeX\ 文档中插入魔方图形变得方便快捷。
本宏包提供的部分命令为方便后续封装的考虑，仅包含在了\env{scope}环境中，用户如果需要直接使用的话，还需要自行添加\env{tikzpicture}环境（或
搭配\cs{tikz}命令来使用），这部分命令我们以\rcs{红色}式样作为标注。
% endregion

% region 命令
\section{命令}

% region \CubeLogo
\begin{function}{\CubeLogo}
  \begin{syntax}
    \cs{CubeLogo}\oarg{放大因子}
  \end{syntax}
绘制一个魔方的标志性图案，默认大小为 1 倍。用法示例如下：
\end{function}
\begin{dispExample*}{sidebyside,lefthand width=0.2\textwidth}
\CubeLogo
\CubeLogo[1.5]
\end{dispExample*}
% endregion

% region \Cube
\begin{function}{\Cube}
  \begin{syntax}
    \cs{Cube}\marg{正面}\marg{右侧面}\marg{顶面}\oarg{单位长度}
  \end{syntax}
绘制一个魔方的3D图，\meta{正面}、\meta{右侧面} 和 \meta{顶面} 分别表示正面、右侧面和顶面，
\meta{单位长度}用来设置\env{tikzpicture}环境的单位长度，默认为 5mm 。
\end{function}
\meta{正面}、\meta{右侧面} 和 \meta{顶面} 参数分别为 9 个或者 5 个字符的字符串，
每个字符代表一个小方块的颜色，其中\meta{正面}和\meta{右侧面}即可以为 9 个字符，
也可以为 5 个字符（表示只绘制如下图所示的 5 个小方块，此时正面剩余的四个位置为
默认的\textcolor{niceblue}{蓝色}，右侧面剩余的位置为默认的
\textcolor{nicered}{红色}）。每个字符根据顺序如下图所示，
依次表示对应位置的小方块颜色：
\begin{dispExample*}{center lower}
\Cube{123456789}{123456789}{123456789}[1cm]
\Cube{12345}{12345}{123456789}[1cm]
\end{dispExample*}

\meta{正面}、\meta{右侧面} 和 \meta{顶面} 参数支持的颜色字符及其对应颜色如下表所示：

{\noindent
\begin{tabular*}{\textwidth}{*{4}{>{\centering\ttfamily}p{1cm}p{1.5cm}}}
  \hline
  字符 & 颜色 & 字符 & 颜色 & 字符 & 颜色 & 字符 & 颜色 \\
  \hline
  - & \textcolor{black!30}{浅灰色} & = & \textcolor{black!65}{深灰色} & g & \textcolor{green!70!black}{绿色} & b & \textcolor{niceblue}{蓝色} \\
  r & \textcolor{nicered}{红色} & w & \fcolorbox{white}{black}{\textcolor{white}{白色}} & o & \textcolor{orange}{橙色} & y & \textcolor{yellow!80!orange!90!black}{黄色} \\
  \hline
\end{tabular*}}

\bigskip
实际使用中，可以根据需要选择合适的字符来表示不同颜色的小方块。例如，下面的例子展示了如何使用不同的颜色字符来绘制魔方的3D图：
\begin{dispExample*}{center lower}
\Cube{----b}{r--r-}{-----b---}
\Cube{---bb-bbb}{-r--rrrrr}{-----b---}
\Cube{rwryroobo}{bbwyyyywg}{gggbbwwow}
\Cube{bbbbbbbbb}{rrrrrrrrr}{yyyyyyyyy}
\end{dispExample*}
从以上可以看出，5字符是9字符的一种子集。5字符的表达方式更简洁，在F2L等场景下更为实用。
% endregion

% region \view
\begin{function}{\view}
  \begin{syntax}
    \cs{view}\marg{视角}
  \end{syntax}
设置魔方的三维视角，\meta{视角} 参数可以取以下值：
\begin{description}[font=\kaishu\ttfamily,left=0pt,itemindent=-28pt,labelsep=2em]
  \item[d] 默认设置，表示正二等轴测投影，左右对称， z轴略短。
  \item[i] 正等轴测投影，即 x、y、z 轴方向的投影长度相等，\cs{CubeLogo}即采用此视角。
  \item[o] 45° 斜投影，即 xy 平面是真实大小，z 轴方向的投影角度为45度，长度为真实长度的 0.5 倍。
  \item[r] 将魔方旋转 180 度后的等轴测投影（可用于演示 CROSS 的情况）。
\end{description}
\end{function}
使用示例如下：
\begin{dispExample*}{center lower}
{\view{r}\Cube{-r--r----}{-b--b----}{-w-www-w-}}
{\view{d}\Cube{---bb}{rr---}{---------}}
{\view{i}\Cube{---bb}{rr---}{yyyyyyyyy}[3mm]}
{\view{o}\Cube{bbbbb}{rrrrr}{yyyyyyyyy}}
\end{dispExample*}
% endregion

% region \CubeNotation
\begin{function}{\CubeNotation}
  \begin{syntax}
    \cs{CubeNotation}\marg{魔方公式符号}\oarg{单位长度}
  \end{syntax}
绘制魔方公式的三维图，\meta{单位长度} 用来设置 \env{tikzpicture} 环境的单位长度，默认为 3mm.
\meta{魔方公式符号} 参数表示魔方公式的各个转动符号。

\cs{CubeNotation} 命令是用几个 coffin 包装起来的完整标记，用法示例：
\end{function}
\begin{dispExample*}{sidebyside,lefthand width=0.3\textwidth}
\CubeNotation{R}
\CubeNotation{b'}
\CubeNotation{M'2}
\CubeNotation{y}
\end{dispExample*}
\cs{CubeNotation} 作为一个用 coffin 封装的命令，它是有固定的高度和宽度的，
因此在排版时可以方便地进行对齐。它们的高度和宽度储存在以下变量中：
\begin{variable}{\gnotationheight,\gnotationwidth,\gcubenotationbox}
  \begin{syntax}
    \cs{gnotationheight} \\
    \cs{gnotationwidth} \\
    \cs{gcubenotationbox}
  \end{syntax}
  这三个变量分别表示魔方公式图的高度、宽度和 box 对象，可以用于自定义排版。
  请注意，视角变换会影响这些变量的值，因此如果要进行全局视角切换，请在
  |\begin{document}| 之前添加 \cs{view} 命令。
\end{variable}
% endregion

% region \FLL
\begin{function}{\FLL}
  \begin{syntax}
    \cs{FLL}\marg{正面}\marg{右侧面}\marg{顶面}\oarg{单位长度}\marg{魔方公式}
  \end{syntax}
  \cs{FLL} 命令用于绘制 F2L 步骤的魔方公式示意图，它是基于\cs{Cube}命令的 coffin 封装，
  \meta{正面}、\meta{右侧面}、\meta{顶面} 和 \oarg{单位长度}参数的含义与
  \cs{Cube}命令相同，\marg{魔方公式} 是可以接收多行输入的魔方公式符号字符串，通常不要超过5行。
  使用示例如下：
\end{function}
\begin{dispExample*}{center lower}
\FLL{----b}{r--r-}{-----b---}{(R'\red{F'}RU)\\(R\blue{U'}R'F)}
\FLL{---bw}{br---}{---------}{(RUR'\blue{U'})\\(RU2R'\blue{U'})\\(RUR')}
\FLL{----w}{b--r-}{-----b---}{(R\blue{U'}R'U)\\(R\blue{U'}R')}
\end{dispExample*}
% endregion

% endregion
\end{document}